\begingroup
\renewcommand{\arraystretch}{1.45}
\fancyhead[L]{\small ZUR THEORIE DER NEUEN TRANSZENDENTEN. III.}

\Needspace{27\baselineskip}
\SourceLeaf{457}
Anwendung auf arithmetisch-geometrische Mittel.
\[
\begin{gathered}
20.\quad a=h(1+2x+2x^4+\mathinner{.\,.})^2\\[6pt]b=h(1-2x+2x^4-\mathinner{.\,.})^2\\[7pt]
a'=h(1+2xx+2x^8+\mathinner{.\,.})^2\\[6pt]b'=h(1-2xx+2x^8-\mathinner{.\,.})^2\\[7pt]
a''=h(1+2x^4+2x^{16}+\mathinner{.\,.})^2\\[6pt]b''=h(1-2x^4+2x^{16}-\mathinner{.\,.})^2\\[7pt]
a'''=h(1+2x^8+2x^{32}+\mathinner{.\,.})^2\\[6pt]b'''=h(1-2x^8+2x^{32}-\mathinner{.\,.})^2\\[5pt]
\qquad\mathrm{usw.}\qquad\mathrm{usw.}
\end{gathered}
\]

\Needspace{36\baselineskip}
\[
\begin{gathered}
c=\surd(aa-bb)\\[8pt]
c'=\surd(a'a'-b'b')\\[7pt]=\tfrac12(a-b)=\dfrac{cc}{4a'}=\dfrac{cc}{4a'},
\\[7pt]\surd\dfrac{c'}{4h}=\dfrac{c}{4}\dfrac{1}{\surd a'h}\\[12pt]
c''=\surd(a''a''-b''b'')\\[7pt]=\tfrac12(a'-b')=\dfrac{c'c'}{4a''}=\dfrac{c^4}{64a'a'a''},
\\[7pt]{}^{4}\!\surd\dfrac{c''}{4h}=\dfrac{c}{4{a'}^{\frac12}{a''}^{\frac14}h^{\frac14}}\\[12pt]
c'''=\surd(a'''a'''-b'''b''')\\[7pt]=\tfrac12(a''-b'')=\dfrac{c''c''}{4a'''}=\dfrac{c^8}{2^{14}{a'}^4{a''}^2a'''},
\\[7pt]{}^{8}\!\surd\dfrac{c'''}{4h}=\dfrac{c}{4{a'}^{\frac12}{a''}^{\frac14}{a'''}^{\frac18}h^{\frac18}}\\[7pt]
\qquad\mathrm{usw.}\\[7pt]\mathrm{usw.}
\end{gathered}
\]

\Needspace{14\baselineskip}
\[
\begin{gathered}21.\quad x^{\frac12}=\frac{c}{4{a'}^{\frac12}{a''}^{\frac14}{a'''}^{\frac18}\mathinner{.\,.}}
\\[8pt]\quad{}=\frac{c}{4a'}\Bigl[\frac{a'}{a''}\Bigr]^{\frac12}
\Bigl[\frac{a''}{a'''}\Bigr]^{\frac14}
\Bigl[\frac{a'''}{a''''}\Bigr]^{\frac18}\ \mathrm{usw.}\end{gathered}
\]

\Needspace{8\baselineskip}
\[
22.\quad x=\frac{a-b}{8a''}\left[\frac{a''}{a'''}\right]^{\frac12}
\left[\frac{a'''}{a''''}\right]^{\frac14}\ \mathrm{usw.}
\]
\Needspace{14\baselineskip}
Ersetzt man in 6 $x$ durch $x^3$ und $y$ durch $-x$, so wird
\[
\begin{gathered}
1-x^4\mathbin{.}1-x^{10}\mathbin{.}1-x^{16}\mathinner{.\,.\,.}
\\[8pt]\quad{}1-x^2\mathbin{.}1-x^8\mathbin{.}1-x^{14}\mathinner{.\,.\,.}\\[8pt]
\qquad=\frac{1}{[x^6]}\{1-xx-x^4+x^{10}+x^{14}-x^{24}-x^{30}+\mathinner{.\,.}\}
\end{gathered}
\]
\Needspace{7\baselineskip}
oder
\[
23.\quad[x]=1-x-xx+x^5+x^7-x^{12}-x^{15}+\mathrm{usw.}
\]
\Needspace{18\baselineskip}
Anderer Beweis.

Man ersetze in 6 $x$ durch $x^3$ und $y$ durch $+x$, so wird

\SourceLeaf{458}
\[
\begin{gathered}
24.\quad 1+x+xx+x^{5}+x^{7}+x^{12}+x^{15}+\mathrm{usw.}
 \\[8pt]\quad{}= [x^{3}]\mathbin{.}1+x\mathbin{.}1+xx\mathbin{.}1+x^{4}\mathbin{.}1+x^{5}\mathinner{.\,.}\\[8pt]
 \\[8pt]\quad{}=\frac{[xx][x^{3}]^{2}}{[x][x^{6}]}
\end{gathered}
\]
\Needspace{14\baselineskip}
Man ersetze in 6 $x$ durch $x^{3}$ und $y$ durch $xxy$, so wird
\[
\begin{gathered}
1+\frac{x}{y}+x^{5}y+\frac{x^{8}}{yy}+x^{16}yy+\mathrm{usw.}\\[8pt]
\qquad=1+\frac{x}{y}\mathbin{.}1+\frac{x^{7}}{y}\mathbin{.}1+\frac{x^{13}}{y}\mathinner{.\,.\,.}
  \\[8pt]\quad{}1+x^{5}y\mathbin{.}1+x^{11}y\mathbin{.}1+x^{17}y\mathinner{.\,.\,.}[x^{6}]
\end{gathered}
\]
\Needspace{14\baselineskip}
oder nach Ersetzung von $x$ durch $x^{3}$
\[
\begin{gathered}
x+\frac{x^{4}}{y}+x^{16}y+\frac{x^{25}}{yy}+x^{49}yy+\mathrm{usw.}\\[8pt]
\qquad=x\mathbin{.}1+\frac{x^{3}}{y}\mathbin{.}1+\frac{x^{21}}{y}\mathbin{.}1+\frac{x^{39}}{y}\mathinner{.\,.\,.}
  \\[8pt]\quad{}1+x^{15}y\mathbin{.}1+x^{33}y\mathbin{.}1+x^{51}y\mathinner{.\,.\,.}[x^{18}]
\end{gathered}
\]
\Needspace{16\baselineskip}
also
\[
\begin{gathered}
25.\quad x+x^{4}+x^{16}+x^{25}+x^{49}+\mathinner{.\,.}\\[8pt]
 =x\mathbin{.}1+x^{3}\mathbin{.}1+x^{15}\mathbin{.}1+x^{21}\\[8pt]\quad{}\mathbin{.}1+x^{33}\mathbin{.}1+x^{39}\mathinner{.\,.\,.}[x^{18}]
 \\[8pt]\quad{}=x\frac{[x^{6}]^{2}[x^{9}][x^{36}]}{[x^{3}][x^{12}][x^{18}]}
\end{gathered}
\]

\Needspace{15\baselineskip}
\[
\begin{gathered}
26.\quad x-x^{4}-x^{16}+x^{25}+x^{49}-\mathrm{usw.}
 \\[8pt]\quad{}=x\mathbin{.}1-x^{3}\mathbin{.}1-x^{15}\mathbin{.}1-x^{21}\mathbin{.}1-x^{33}\mathinner{.\,.}[x^{18}]\\[8pt]
 \\[8pt]\quad{}=x\frac{[x^{3}][x^{18}]^{2}}{[x^{6}][x^{9}]}
\end{gathered}
\]
\Needspace{28\baselineskip}
Da nun
\[
\begin{gathered}\frac{x\mathbin{.}1-x^{3}\mathbin{.}1-x^{15}\mathbin{.}1-x^{21}\mathbin{.}1-x^{33}\mathinner{.\,.\,.}[x^{18}]}
     {x\mathbin{.}1-x^{3}\mathbin{.}1-x^{9}\mathbin{.}1-x^{15}\mathbin{.}1-x^{21}\mathinner{.\,.\,.}}
 \\[8pt]\quad{}=\frac{[x^{18}]^{2}}{[x^{9}]}
\\[8pt]\quad{} =1+x^{9}+x^{27}+x^{54}+\mathinner{.\,.}\end{gathered}
\]
so ist
\[
\begin{gathered}
27.\quad x-x^{4}-x^{16}+x^{25}+x^{29}-\mathinner{.\,.}\\[8pt]
 =1-x^{3}\mathbin{.}1-x^{9}\mathbin{.}1-x^{15}\mathinner{.\,.\,.}\frac{1}{x^{\frac18}}
 \\[8pt]\quad{}\{x^{\frac98}+x^{\frac{81}{8}}+x^{\frac{225}{8}}+x^{\frac{441}{8}}+\mathrm{usw.}\}
\end{gathered}
\]
\Needspace{27\baselineskip}
Ferner folgt aus 24, wenn man $x$ durch $x^{3}$ ersetzt, da
\[
\begin{gathered}\frac{1-2x^{9}+2x^{36}-2x^{81}+\mathinner{.\,.\,.}}
     {1-x^{3}\mathbin{.}1-x^{9}\mathbin{.}1-x^{15}\mathinner{.\,.\,.}}
 \\[8pt]\quad{}=\frac{[x^{6}][x^{9}]^{2}}{[x^{3}][x^{18}]}
\\[8pt]\quad{} =1+x^{3}+x^{6}+x^{15}+x^{21}+\mathinner{.\,.}\end{gathered}
\]
\[
\begin{gathered}28.\quad 1-2x^{9}+2x^{36}-2x^{81}+\mathinner{.\,.}
\\[8pt]\quad{} =1-x^{3}\mathbin{.}1-x^{9}\mathbin{.}1-x^{15}\mathinner{.\,.\,.}\frac{1}{x^{\frac18}}
 \\[8pt]\quad{}\{x^{\frac18}+x^{\frac{25}{8}}+x^{\frac{49}{8}}+\mathrm{usw.}\}\end{gathered}
\]
\Needspace{16\baselineskip}
also aus der Verbindung von 27 und 28
\SourceLeaf{459}
\[
\begin{gathered}
29.\quad 1-2x+2x^{4}-2x^{9}+2x^{16}-\mathinner{.\,.}\\[8pt]
 =1-x^{3}\mathbin{.}1-x^{9}\mathbin{.}1-x^{15}\mathinner{.\,.\,.}\frac{1}{x^{\frac18}}
 \\[8pt]\quad{}\{x^{\frac18}-2x^{\frac98}+x^{\frac{25}{8}}+x^{\frac{49}{8}}-2x^{\frac{81}{8}}-\mathrm{usw.}\}
\end{gathered}
\]
\Needspace{24\baselineskip}
Ferner folgt aus 23
\[
x^{\frac18}-x^{\frac{25}{8}}-x^{\frac{49}{8}}+x^{\frac{121}{8}}+\mathinner{.\,.}
 =x^{\frac18}[x^{3}]
\]
also
\[
\begin{gathered}
30.\quad 1-2x^{3}+2x^{12}-2x^{27}+2x^{48}-\mathinner{.\,.}\\[8pt]
 =1-x^{3}\mathbin{.}1-x^{9}\mathbin{.}1-x^{15}\mathinner{.\,.\,.}\frac{1}{x^{\frac18}}
 \\[8pt]\quad{}\{x^{\frac18}-x^{\frac{25}{8}}-x^{\frac{49}{8}}+x^{\frac{121}{8}}+x^{\frac{169}{8}}-\mathinner{.\,.}\}
\end{gathered}
\]
\Needspace{20\baselineskip}
Aus der Addition von 29 und 30
\[
\begin{gathered}
31.\quad \{1-2x+2x^{4}-2x^{9}+\mathinner{.\,.}\}\\[8pt]\quad{}+\{1-2x^{3}+2x^{12}-2x^{27}+\mathinner{.\,.}\}\\[8pt]
 =1-x^{3}\mathbin{.}1-x^{9}\mathbin{.}1-x^{15}\mathinner{.\,.\,.}\frac{2}{x^{\frac18}}
 \\[8pt]\quad{}\{x^{\frac18}-x^{\frac98}-x^{\frac{81}{8}}+x^{\frac{121}{8}}+x^{\frac{169}{8}}-\mathinner{.\,.}\}
\end{gathered}
\]
\Needspace{20\baselineskip}
Aus der Addition von 28 und 30
\[
\begin{gathered}
32.\quad \{1-2x^{3}+2x^{12}-2x^{27}+\mathinner{.\,.}\}\\[8pt]\quad{}+\{1-2x^{9}+2x^{36}-2x^{81}+\mathinner{.\,.}\}\\[8pt]
 =1-x^{3}\mathbin{.}1-x^{9}\mathbin{.}1-x^{15}\mathinner{.\,.\,.}\frac{2}{x^{\frac18}}
 \\[8pt]\quad{}\{x^{\frac18}+x^{\frac{121}{8}}+x^{\frac{169}{8}}+x^{\frac{529}{8}}+\mathinner{.\,.}\}
\end{gathered}
\]
\clearpage
\endgroup
